% 第 IX 讲，第14节至参考文献；原 PDF 515--532。
\SourcePage{515}
\SGAsection{14}{与超越理论的联系：刚性解析情形。典范同态 $M_K\longrightarrow A_K^{\natural0}$}
\textbf{14.1。}这里假定读者熟悉\Term{rigid-analytic-space}{刚性解析空间}理论的基础；可参阅 [12]、[13]、[34]。

设 $S$ 为\Term{complete}{完备}\Term{trait}{离散赋值环谱}，$A_K$ 为其\Term{fraction-field}{分式域}上的\Term{abelian-scheme}{阿贝尔概形}。8.1 中已从\Term{semistable}{半稳定}情形的\Term{raynaud-extension}{Raynaud 扩张}出发，经\Term{descent}{下降}构造了 $G_K=A_K^{\natural0}$ 的\Term{extension}{扩张} (8.1.1)：
\begin{equation}\tag{14.1.1}
0\longrightarrow T_K\longrightarrow G_K\longrightarrow B_K\longrightarrow0,
\end{equation}
即\Term{abelian-scheme}{阿贝尔概形} $B_K$ 被\Term{torus}{环面} $T_K$ 的\Term{extension}{扩张}。考察\Term{dual-abelian-scheme}{对偶阿贝尔概形} $A_K'$，同样得到 $G_K'=A_K'^{\natural0}$ 的\Term{extension}{扩张}
\begin{equation}\tag{14.1.2}
0\longrightarrow T_K'\longrightarrow G_K'\longrightarrow B_K'\longrightarrow0,
\end{equation}
其中 $B_K'$ 还可与 $B_K$ 的\Term{dual-abelian-scheme}{对偶阿贝尔概形}相认同（8.2）。与 (8.1.8) 一样，分别以 $M_K,M_K'$ 表示 $K$ 上与 $T_K',T_K$ \Term{dual}{对偶}的\Term{lattice}{格}：
\begin{equation}\tag{14.1.3}
M_K=D(T_K'),\qquad M_K'=D(T_K).
\end{equation}
\SourcePage{516}
因此，\Term{extension}{扩张}结构 (14.1.2) 和 (14.1.1) 分别\Term{solvable}{可解}释为（第 VIII 讲 3.7）\Term{equivalence}{等价}于给定\Term{homomorphism}{同态}
\begin{equation}\tag{14.1.4}
M_K\xrightarrow{p_K}B_K,\qquad M_K'\xrightarrow{p_K'}B_K'.
\end{equation}
设 $W_K^\natural$ 为定义这些\Term{abelian-scheme}{阿贝尔概形}之间\Term{canonical}{典范}\Term{duality}{对偶性}的、$(B_K,B_K')$ 被 $\mathbb G_{mK}$ 的\Term{biextension}{双扩张}（8.2），并设 $E_K$ 为 $W_K^\natural$ 沿\Term{homomorphism}{同态} (14.1.4) 的\Term{inverse-image}{逆像}，即 $(M_K,M_K')$ 被 $\mathbb G_{mK}$ 的\Term{biextension}{双扩张}。注意，立即可以验证，$E_K$ 的\Term{trivialization}{平凡化}与\Term{lifting}{提升} $p_K:M_K\to B_K$ 的\Term{homomorphism}{同态}
\begin{equation}\tag{14.1.5}
i:M_K\longrightarrow G_K
\end{equation}
之间存在\Term{canonical}{典范}的一一对应；\Term{symmetric}{对称}地，它们也与\Term{lifting}{提升} $p_K':M_K'\to B_K'$ 的\Term{homomorphism}{同态}
\begin{equation}\tag{14.1.6}
i':M_K'\longrightarrow G_K'
\end{equation}
\Term{canonical}{典范}地一一对应。

现在，以上标 $\mathrm{an}$ 表示与所考察的\Term{algebraic-group}{代数群}关联的 $K$ 上的\Term{rigid-analytic-group}{刚性解析群} [13]。文献 [24] 定义了一个\Term{canonical}{典范}\Term{exact-sequence}{正合列}
\begin{equation}\tag{14.1.7}
0\longrightarrow M_K^{\mathrm{an}}\xrightarrow{i^{\mathrm{an}}}
G_K^{\mathrm{an}}\xrightarrow e A_K^{\mathrm{an}}\longrightarrow0,
\end{equation}
其中\Term{homomorphism}{同态} $i^{\mathrm{an}}$ 也\Term{solvable}{可解}释为通常\Term{group-scheme}{群概形}之间的\Term{homomorphism}{同态} $i:M_K\to G_K$，它\Term{lifting}{提升}\Term{homomorphism}{同态} $p$ (14.1.4)，也就是说，它具有 (14.1.5) 所考察的\Term{form}{形式}。这个\Term{exact-sequence}{正合列}是 (13.3.4) 的\Term{rigid-analytic}{刚性解析}对应物。它以\Term{functorial}{函子性}方式依赖于 $A_K$，
\SourcePage{517}
而且它的形成与\Term{trait}{离散赋值环谱}的任意\Term{base-change}{基变换} $S'\to S$ 相容。后一个事实还表明，要定义和研究 (14.1.7)，可经\Term{finite-descent}{有限下降}归结到 $A_K$ 具有\Term{semistable-reduction}{半稳定约化}的情形。这时，\Term{group-scheme}{群概形} $T_K,G_K,B_K,M_K,\ldots$ 是 $S$ 上\Term{group-scheme}{群概形} $T,G,B,M,\ldots$ 的\Term{generic-fiber}{泛纤维}，而 $G$ 是\Term{abelian-scheme}{阿贝尔概形} $B$ 被\Term{torus}{环面} $T$ 的\Term{extension}{扩张}（第 7 节）。于是可以考察 $G$ 沿其\Term{special-fiber}{特殊纤维}的\Term{formal-completion}{形式完备化} $\widehat G$；由构造，有\Term{canonical}{典范}\Term{isomorphism}{同构}（7.2.6）
\begin{equation}\tag{14.1.8}
\widehat G\xrightarrow\sim\widehat{A^0},
\end{equation}
从而取“\Term{generic-fiber}{泛纤维}”（它们是 $K$ 上的\Term{rigid-analytic-group}{刚性解析群}），得到\Term{canonical}{典范}\Term{isomorphism}{同构}
\begin{equation}\tag{14.1.9}
\widehat G^{\mathrm{an}}\xrightarrow\sim\widehat{A^0}^{\mathrm{an}}.
\end{equation}
另一方面，有\Term{canonical}{典范}\Term{homomorphism}{同态}（它们是\Term{rigid-analytic}{刚性解析}的\Term{open-immersion}{开浸入}，但\Term{generic}{一般}不是\Term{isomorphism}{同构}）
\begin{equation}\tag{14.1.10}
\widehat G^{\mathrm{an}}\lhook\joinrel\longrightarrow G_K^{\mathrm{an}},
\qquad \widehat{A^0}^{\mathrm{an}}\lhook\joinrel\longrightarrow A_K^{\mathrm{an}}.
\end{equation}
这样，(14.1.7) 中的\Term{homomorphism}{同态} $e$ 由下图的\Term{commutative-diagram}{交换性}唯一刻画：
\begin{equation}\tag{14.1.11}
\begin{tikzcd}
\widehat G^{\mathrm{an}}\arrow[r,"\sim"]\arrow[d]&\widehat{A^0}^{\mathrm{an}}\arrow[d]\\
G_K^{\mathrm{an}}\arrow[r,"e"]&A_K^{\mathrm{an}},
\end{tikzcd}
\end{equation}
其中未标名的\Term{arrow}{箭头}来自 (14.1.9) 和 (14.1.10)。

\SourcePage{518}
\textbf{注记 14.1.12。}a) 注意，使上图\Term{commutative-diagram}{交换}的 $e$ 的唯一性是直接的，因为\Term{rigid-analytic-space}{刚性解析空间} $G_K^{\mathrm{an}}$ \Term{connected}{连通}。但 $e$ 的存在性绝非平凡；关于这一点，我们仅请读者参阅前引 M. Raynaud 的短文。另一方面，该文中给出的\Term{homomorphism}{同态} $i^{\mathrm{an}}$ 的定义也具有\Term{transcendental}{超越}性质，尽管它\Term{solvable}{可解}释为一个本质上\Term{algebra}{代数}性的\Term{object}{对象} (14.1.5)（一旦已经构造出 $G$ 本身，因而构造出 $G_K$；这一步由第 7 节的\Term{formal-geometry}{形式几何}构造完成）。我不知道它有任何纯\Term{algebraic-geometry}{代数几何}的描述，但下文 (14.4) 将提出一个“算术”刻画，它可通过归结到环 $\mathbb Z$ 上的\Term{finite-type}{有限型}情形得到。因此，从\Term{algebraic-geometry}{代数几何}的观点看来，\Term{homomorphism}{同态} (14.1.5) 似乎是一种比前面各节遇到的结构更为隐蔽的结构。

b) 如上所述，\Term{homomorphism}{同态} $i$ \Term{canonical}{典范}地对应于 $(M_K,M_K')$ 的\Term{biextension}{双扩张} $E_K$ 的某个确定的\Term{trivialization}{平凡化}，也对应于\Term{lifting}{提升} $p_K'$ 的一个确定的\Term{homomorphism}{同态} (14.1.6)。可以证明，后者正是与 $A_K$ 的\Term{dual-abelian-scheme}{对偶阿贝尔概形} $A_K'$ 关联、类似于 (14.1.7) 的\Term{exact-sequence}{正合列}
\begin{equation}\tag{14.1.13}
0\longrightarrow M_K'^{\mathrm{an}}\xrightarrow{i'^{\mathrm{an}}}
G_K'^{\mathrm{an}}\longrightarrow A_K'^{\mathrm{an}}\longrightarrow0
\end{equation}
中出现的\Term{homomorphism}{同态}（除非应当取其相反数？？？）。

\SourcePage{519}
\textbf{14.2。\Term{homomorphism}{同态} $i:M_K\to G_K$ 与\Term{monodromy}{单值性}的关系。}
设 $n$ 是正整数。在 $K$ 上\Term{rigid-analytic-space}{刚性解析空间}的\Term{site}{景}上，利用 $\operatorname{Ext}_{\mathbb Z}^i(\mathbb Z/n\mathbb Z,-)$ 的\Term{exact-sequence}{正合列}，可由\Term{exact-sequence}{正合列} (14.1.7) 得到一个\Term{canonical}{典范}\Term{exact-sequence}{正合列}（注意 $M_K^{\mathrm{an}}$ \Term{torsion-free}{无挠}，且 $G_K^{\mathrm{an}}$ 是 $n$-\Term{divisible}{可除}的）：
\[
0\longrightarrow{}_nG_K^{\mathrm{an}}\longrightarrow{}_nA_K^{\mathrm{an}}
\longrightarrow M_{K,n}^{\mathrm{an}}\longrightarrow0,
\qquad M_{K,n}\mathrel{\mathop:}=M_K/nM_K.
\]
由于这\Term{exact-sequence}{正合列}中的所有项均来自 $K$ 上的\Term{finite-group-scheme}{有限群概形}，可将它改写为来自\Term{exact-sequence}{正合列}
\begin{equation}\tag{14.2.1}
0\longrightarrow{}_nG_K\longrightarrow{}_nA_K\longrightarrow M_{K,n}\longrightarrow0.
\end{equation}
当 $n$ 变化时，这些\Term{exact-sequence}{正合列}具有应有的相容性，因而对某个\Term{prime-number}{素数} $\ell$ 的各次幂取\Term{limit}{极限}，便得到类似的\Term{exact-sequence}{正合列}
\begin{equation}\tag{14.2.2}
0\longrightarrow T_\ell(G_K)\longrightarrow T_\ell(A_K)
\longrightarrow M_K\otimes_{\mathbb Z}\mathbb Z_\ell\longrightarrow0.
\end{equation}
而利用 (8.1.2) 的第一个公式和 (8.1.7)，我们已由纯\Term{algebra}{代数}方法得到这样一个\Term{exact-sequence}{正合列}，因为它们给出
\begin{equation}\tag{14.2.3}
\begin{gathered}
T_\ell(G_K)\xrightarrow\sim T_\ell(A_K)^{ef}\subset T_\ell(A_K),\\
T_\ell(A_K)/T_\ell(A_K)^{ef}\xrightarrow\sim M_K\otimes_{\mathbb Z}\mathbb Z_\ell.
\end{gathered}
\end{equation}
现在，我断言：来自\Term{algebra}{代数}理论的\Term{homomorphism}{同态} (14.2.3)，正是从\Term{transcendental}{超越}理论给出的\Term{exact-sequence}{正合列} (14.2.2) \Term{derive}{导出}的那些\Term{homomorphism}{同态}。对于与 (14.2.2) 的第一个\Term{homomorphism}{同态}对应的第一个\Term{homomorphism}{同态}，先归结到\Term{semistable}{半稳定}情形，
\SourcePage{520}
我们的断言便由 (14.1.11) 的\Term{commutative-diagram}{交换性}立即得出：其中的竖\Term{arrow}{箭头}在 $T_\ell$ 上定义\Term{isomorphism}{同构}，而对横\Term{arrow}{箭头}应用 $T_\ell$，分别得到 (14.2.3) 和 (14.2.2) 中考察的\Term{arrow}{箭头}。至于由此得到的 $T_\ell(A_K)/T_\ell(A_K)^{ef}$ 与 $M_K\otimes\mathbb Z_\ell$ 之间两个\Term{isomorphism}{同构}的上述相容性，我们将予以承认。还应注意，这个相容性（对某个给定的 $\ell$）已经足以用\Term{transcendental}{超越}方法唯一刻画 (14.1.7) 中出现的\Term{homomorphism}{同态} (14.1.5)，即把它解释为 $M_K^{\mathrm{an}}$ 与 $\operatorname{Ker}e$ 之间的\Term{isomorphism}{同构}。也可以说，(14.1.7) 的\Term{homomorphism}{同态} $e$ 使我们得以把各个 $T_\ell(A_K)/T_\ell(A_K)^{ef}$ 与 $K$ 上同一个\Term{lattice}{格} $R_K=\operatorname{Ker}e$ 的 $\ell$-进化 $R_K\otimes_{\mathbb Z}\mathbb Z_\ell$ 相认同；上述\Term{pro-group}{pro-$\ell$-群}系统上的这个“整结构”，在显然的意义下与公式 (8.1.7)（\Term{orthogonality}{正交}\Term{theorem}{定理} 5.2）所给的整结构相容，因而\Term{canonical}{典范}地定义一个\Term{isomorphism}{同构}
\[
i:M_K\xrightarrow\sim\operatorname{Ker}e.
\]
由于\Term{homomorphism}{同态} $i$ (14.1.5) 能恢复 $T_\ell(A_K)$ 的通常\Term{extension}{扩张}结构 (14.2.2)，再利用由 (14.1.1) 得到的\Term{exact-sequence}{正合列}
\[
0\longrightarrow\underset{\substack{\Vert\\M_K'\otimes\mathbb Z_\ell(1)}}{T_\ell(T_K)}
\longrightarrow T_\ell(G_K)\longrightarrow T_\ell(B_K)\longrightarrow0,
\]
也能恢复 9.6 所描述的\Term{blended-extension}{混合扩张}结构，
\SourcePage{521}
所以在\Term{semistable}{半稳定}情形，它同样能恢复\Term{monodromy-pairing}{单值性配对} (9.1.2)
\begin{equation}\tag{14.2.4}
u_\ell:M\otimes M'\longrightarrow\mathbb Z_\ell,
\end{equation}
因为这些\Term{pairing}{配对}在 9.6 中正是用\Term{blended-extension}{混合扩张}定义的。这里仅说明将 9.6 的构造转译后得到的结果。为此，注意在\Term{semistable}{半稳定}情形，$p_K,p_K'$ 来自\Term{homomorphism}{同态}
\[
p:M\longrightarrow B,\qquad M'\longrightarrow B',
\]
其中 $B,B'$ 为 $S$ 上互相\Term{dual}{对偶}的两个\Term{abelian-scheme}{阿贝尔概形}。于是沿 $(p,p')$ 取\Term{inverse-image}{逆像}，得到 $(M,M')$ 被 $\mathbb G_{mS}$ 的\Term{biextension}{双扩张} $E$，其\Term{generic-fiber}{泛纤维}就是 4.1 中已考察的\Term{biextension}{双扩张} $E_K$。现在为简便起见，假定 $M,M'$ 为\Term{constant}{常值}，分别取值于 $M,M'$。对每一对 $(m,m')\in M\times M'$，$E_{m,m'}$ 是 $\mathbb G_{mS}$ 下的\Term{torsor}{挠子}，因而对应于 $S$ 的环 $V$ 上的一个\Term{invertible-module}{可逆模} $L_{m,m'}$。所以它有一个基 $u_0$，这个基在乘以 $V$ 的\Term{unit}{单位}的意义下确定。由此得到一个\Term{canonical}{典范}的“\Term{valuation}{赋值}”映射
\[
\underline v:E_{m,m'}(K)\longrightarrow\mathbb Z,
\]
即对 $L_{m,m'}\otimes_VK$ 的任一基 $u$，将它写成 $u=\lambda u_0$（$\lambda\in K^*$），并赋予它元素
\[
\underline v(u)=v(\lambda)\in\mathbb Z,
\]
其中右端的 $v$ 表示与 $V$ 关联的\Term{valuation}{赋值}（按通常方式规范化，使其\Term{value-group}{值群}为 $\mathbb Z$）；
\SourcePage{522}
所得整数显然不依赖于 $u_0$ 的选择。若 $\varphi$ 是 $E$ 的一个\Term{trivialization}{平凡化}，便可考察相应的映射
\begin{equation}\tag{14.2.5}
(m,m')\longmapsto\underline v(\varphi(m,m')):M\times M'\longrightarrow\mathbb Z.
\end{equation}
这样，可以得到：\Term{monodromy-homomorphism}{单值性同态} (14.2.4) 取值于 $\mathbb Z$，而且它正是\Term{homomorphism}{同态} (14.2.5)，其中 $\varphi$ 是与\Term{homomorphism}{同态} $i$ (14.1.5) 关联的 $E_K$ 的\Term{splitting}{分裂}。

\textbf{注记 14.2.6。}特别地，由此（通过\Term{rigid-analytic}{刚性解析}的\Term{transcendental}{超越}理论）重新得到\Term{monodromy-pairing}{单值性配对}的\Term{integrality}{整性}与对 $\ell$ 的独立性（10.4 a)）；12.5 已用另一种（同样具有\Term{transcendental}{超越}性质的）方法证明了这一点。据 M. Raynaud 所言，\Term{rigid-analytic}{刚性解析}理论也能重新得到正性断言 10.4 b)（它同样包含在 12.5 中）。

\textbf{注记 14.2.7。}把 $M_K\otimes_{\mathbb Z}\mathbb Q$ 看成 $M_K\otimes_{\mathbb Z}\mathbb Q/\mathbb Z$ 被 $M_K$ 的\Term{extension}{扩张}，再利用\Term{homomorphism}{同态} (14.1.5) $i:M_K\to G_K$，便得到一个\Term{extension}{扩张} $A_K^\flat$：
\begin{equation}\tag{14.2.8}
0\longrightarrow G_K\longrightarrow A_K^\flat\longrightarrow
M_K\otimes_{\mathbb Z}\mathbb Q/\mathbb Z\longrightarrow0.
\end{equation}
还应注意，由于 $G_K$ 是\Term{divisible-group}{可除群}，给定这样一个\Term{extension}{扩张}，\Term{equivalence}{等价}于给定一族形如 (14.2.2) 的\Term{extension}{扩张}（即对 (14.2.8) 应用 $T_\ell$ 得到的那些\Term{extension}{扩张}）。而且，由于 $i$ 是\Term{monomorphism}{单态射}，定义 (14.2.8) 的这些\Term{extension}{扩张}正是从\Term{exact-sequence}{正合列} (14.2.1) 得到的\Term{extension}{扩张} (14.2.2)。
\SourcePage{523}
由于后者由纯\Term{algebra}{代数}方法定义，\Term{extension}{扩张} (14.2.8) 也是如此。事实上，刚给出的刻画立即表明，考虑到 11.11 将 (8.2.3) 所引入的、同样记作 $A_K^\flat$ 的群的\Term{connected-component-group}{连通分支群}与 $M_K\otimes\mathbb Q/\mathbb Z$ 相认同的\Term{isomorphism}{同构}，这个\Term{extension}{扩张}与该群\Term{canonical}{典范}\Term{isomorphism}{同构}。

\textbf{14.3。应用于 $i:M_K\to G_K$ 的算术刻画。}
我们所要给出的刻画基于如下观察：对每个\Term{prime-number}{素数} $\ell$，从 $i$ 得到的 $M_K\otimes_{\mathbb Z}\mathbb Z_\ell$ 被 $T_\ell(G_K)$ 的\Term{extension}{扩张} (14.2.2)，是由纯\Term{algebra}{代数}方法已知的；特别地，我们知道 $i$ 在\Term{canonical}{典范}\Term{homomorphism}{同态}
\begin{equation}\tag{14.3.1}
\operatorname{Hom}(M_K,G_K)\longrightarrow
\prod_\ell\operatorname{Ext}^1(M_K\otimes_{\mathbb Z}\mathbb Z_\ell,T_\ell(G_K))
\end{equation}
下的像。在 $M_K$ 为\Term{constant}{常值}、取值于 $M$ 的情形（经\Term{descent}{下降}总可归结到这种情形），上述\Term{homomorphism}{同态}\Term{solvable}{可解}释为\Term{homomorphism}{同态}
\begin{equation}\tag{14.3.2}
\operatorname{Hom}(M,G_K(K))\longrightarrow
\prod_\ell\operatorname{Hom}(M,H^1(K,T_\ell(G_K))),
\end{equation}
它来自\Term{canonical}{典范}\Term{homomorphism}{同态}
\begin{equation}\tag{14.3.3}
G_K(K)\longrightarrow\prod_\ell H^1(K,T_\ell(G_K)),
\end{equation}
而后者来自 Kummer \Term{exact-sequence}{正合列}中的\Term{boundary-homomorphism}{上边缘同态}：
\[
0\longrightarrow{}_nG_K\longrightarrow G_K\xrightarrow n G_K\longrightarrow0.
\]
\SourcePage{524}
（其实，\Term{homomorphism}{同态} (14.3.1) 在任何情形下，都\Term{solvable}{可解}释为与 (14.3.3) 类似的\Term{boundary-homomorphism}{上边缘同态}，只是针对\Term{twisting}{扭曲}的群 $M_K^*\otimes G_K=\underline{\operatorname{Hom}}(M_K,G_K)$。）(14.3.3) 的\Term{kernel}{核}显然由 $G_K(K)$ 中对每个整数 $n\geq1$ 都能被 $n$ 整除的元素组成；我们称这样的元素为\Term{infinitely-divisible}{无限可除}元素。因此，这已经在模去 $\operatorname{Hom}(M,g)$ 中的一个元素的意义下确定了 $i$，其中 $g$ 为 $G_K(K)$ 中所有\Term{infinitely-divisible}{无限可除}元素组成的\Term{subgroup}{子群}：
\[
g=\bigcap_{n\geq1}nG_K(K).
\]
现在有如下结果。
\begin{SGAresult}{\Term{lemma}{引理}}{14.3.4}
设 $V$ 为\Term{discrete-valuation-ring}{离散赋值环}，其\Term{residue-field}{剩余域} $k$ 的\Term{characteristic}{特征}为 $p>0$，且在\Term{prime-field}{素域}上为\Term{finite-type}{有限型}；$G_K$ 为 $K$ 上的\Term{finite-type}{有限型}\Term{group-scheme}{群概形}，并且当 $\operatorname{char}K=0$ 时其\Term{unipotent-rank}{幺幂秩}为零。那么，$G_K(K)$ 中\Term{infinitely-divisible}{无限可除}元素组成的群 $g$ 为零；也就是说，（若 $G_K$ \Term{divisible}{可除}，则）\Term{homomorphism}{同态} (14.3.3) 为\Term{injective}{单射}。
\end{SGAresult}
\begin{proof}
将 $K$ 换成一个\Term{finite-extension}{有限扩张} $K'$，并将 $V$ 换成它在 $K'$ 中的\Term{normalization}{正规化}后，可以假定 $G_K$ 有一条\Term{composition-series}{合成列}，其各相邻\Term{quotient}{商}属于下列类型之一：\Term{finite}{有限}\Term{constant-group}{常值群}、\Term{abelian-variety}{阿贝尔簇}、$\mathbb G_{mK}$、$\mathbb G_{aK}$。于是立即归结到 $G_K$ 本身属于上述某一类型的情形。\Term{finite}{有限}\Term{constant-group}{常值群}或群 $\mathbb G_{aK}$ 的情形是平凡的。在\Term{abelian-scheme}{阿贝尔概形}的情形，可将它看成其\Term{neron-model}{Néron 模型} $G$ 的\Term{generic-fiber}{泛纤维}，所以这时
\SourcePage{525}
只须证明：若 $G$ 为 $S$ 上的\Term{finite-type}{有限型}\Term{scheme}{概形}，则 $G(S)$ 中\Term{infinitely-divisible}{无限可除}点组成的\Term{subgroup}{子群}为零。群 $G(k)$ 确实具有这个性质；这立即来自对 $G_0$ 的通常\Term{resolution}{分解}和 Mordell--Weil \Term{theorem}{定理} [16]（该\Term{theorem}{定理}说，若 $G_0$ 为\Term{finite-type}{有限型}域 $k$ 上的\Term{abelian-scheme}{阿贝尔概形}，则群 $G_0(k)$ 是一个\Term{finitely-generated-module}{有限生成模}，系数环为 $\mathbb Z$）。故 $g$ 包含在 $G(S)\to G(k)$ 的\Term{kernel}{核} $G(S)^{(1)}$ 中。由于 $G(S)^{(1)}\to G(V/\mathfrak m^{n+1})$ 的\Term{kernel}{核}\Term{decomposable}{可分解}为与 $k$ 上\Term{vector-group}{向量群}\Term{isomorphism}{同构}的群（这里可假设 $S$ \Term{complete}{完备}），因而这些群都被 $p$ \Term{annihilate}{零化}，所以无论 $n$ 为何，$g$ 都包含在上述\Term{homomorphism}{同态}的\Term{kernel}{核}中，从而 $g$ 为零。$G_K=\mathbb G_{mK}$ 的情形可归结到前一情形：此时 $G_K(K)\xrightarrow\sim K^*$ 是 $\mathbb Z$（其\Term{infinitely-divisible}{无限可除}元素只有零）被 $V^*=\mathbb G_m(S)$ 的\Term{extension}{扩张}，而由刚才的论证，后者也满足同样的条件。
\end{proof}

\textbf{注记 14.3.4.1。}注意，在 14.3.4 中，要求对所有整数 $\geq1$ \Term{divisible}{可除}，而不只是对某个固定\Term{prime-number}{素数} $\ell$ 的各次幂\Term{divisible}{可除}，是至关重要的。换言之，必须考察\Term{homomorphism}{同态} (14.3.3) 的\Term{kernel}{核}，不能在右端只取一个因子 $H^1(K,T_\ell(G_K))$。事实上，若 $G_K$ 来自 $S$ 上的\Term{smooth}{光滑}\Term{group-scheme}{群概形} $G$，且 $S$ 为\Term{henselian}{Hensel}的，则 $G(S)\to G(k)$ 的\Term{kernel}{核}中每个元素对所有 $\ell\ne p$ 都无限 $\ell$-\Term{divisible}{可除}。证明 14.3.4 的论证表明，$G_K(K)$ 中
\SourcePage{526}
无限 $p$-\Term{divisible}{可除}元素组成的\Term{subgroup}{子群}（即当 $G_K$ 为 $p$-\Term{divisible}{可除}时，$G_K(K)\to H^1(K,T_p(G_K))$ 的\Term{kernel}{核}）是\Term{finite}{有限}群。

\textbf{14.3.5。}上述讨论表明，当\Term{residue-field}{剩余域} $k$ 的\Term{characteristic}{特征}为 $p>0$，且在绝对意义下为\Term{finite-type}{有限型}时，知道\Term{homomorphism}{同态} $i$ (14.1.5) 在 (14.3.1) 下的像，即知道各个\Term{extension}{扩张} (14.2.2)，就完全知道了 $i$。为了从中\Term{derive}{导出}不对 $k$ 作\Term{restriction}{限制}的\Term{generic}{一般}情形下刻画 $i$ 的原则，注意第 7 节中已在远比“$S$ 为\Term{complete}{完备}\Term{trait}{离散赋值环谱}，并在其\Term{fraction-field}{分式域}上给定具有\Term{semistable-reduction}{半稳定约化}的\Term{abelian-variety}{阿贝尔簇}”更\Term{generic}{一般}的条件下，构造了 $S$ 上的\Term{scheme}{概形} $G$。具体地，现在更\Term{generic}{一般}地以 $S$ 表示一个\Term{normal}{正规}\Term{noetherian}{诺特}环 $V$ 的\Term{spectrum}{谱}，$V$ 对某个 $\mathfrak m$-\Term{adic-topology}{进拓扑}\Term{separated}{分离}且\Term{complete}{完备}；并考察 $S$ 上一个具有\Term{connected}{连通}\Term{fiber}{纤维}的\Term{smooth}{光滑}\Term{quasi-projective}{拟射影}\Term{group-scheme}{群概形} $A$。假定
$A_0=A\times_SS_0$（$S_0=\operatorname{Spec}(V/\mathfrak m)$）是\Term{abelian-scheme}{阿贝尔概形} $B_0$ 被\Term{isotrivial-torus}{等平凡环面} $T_0$ 的\Term{extension}{扩张}，并且 $A|U$ 为\Term{abelian-scheme}{阿贝尔概形}，其中 $U\subset S-S_0$ 是\Term{nonempty}{非空}\Term{open-subscheme}{开子概形}。于是，7.2.2 中已构造了 $S$ 上的\Term{scheme}{概形} $G$，它是\Term{abelian-scheme}{阿贝尔概形} $B$ 被\Term{torus}{环面} $T$ 的\Term{extension}{扩张}；$G$ 在唯一\Term{isomorphism}{同构}意义下由沿 $G_0,A_0$ 的\Term{formal-completion}{形式完备化}上的\Term{isomorphism}{同构} (14.1.8)
\begin{equation}\tag{14.3.5.1}
\widehat G\xrightarrow\sim\widehat A
\end{equation}
刻画。由此，对每个\Term{prime-number}{素数} $\ell$，得到一个\Term{isomorphism}{同构}（7.3.5）
\SourcePage{527}
\begin{equation}\tag{14.3.5.2}
T_\ell(G)\xrightarrow\sim T_\ell(A)^f,
\end{equation}
从而\Term{restriction}{限制}到 $U$，得到 Barsotti--Tate \Term{pro-group}{pro-$\ell$-群}的\Term{extension}{扩张}结构
\begin{equation}\tag{14.3.5.3}
0\longrightarrow\underset{\substack{\Vert\\T_\ell(A_U)^f}}{T_\ell(G_U)}
\longrightarrow T_\ell(A_U)\longrightarrow T_\ell(A_U)/T_\ell(A_U)^f\longrightarrow0.
\end{equation}
假定给定 $S$ 上的第二个群 $G'$，满足与 $A$ 相同的假设，以及 $(A,A')$ 被 $\mathbb G_m$ 的\Term{biextension}{双扩张} $W$，它在 $U$ 上诱导\Term{abelian-scheme}{阿贝尔概形} $A_U,A_U'$ 之间的\Term{duality}{对偶性}。（事实上，由第 VIII 讲 7.1，给定 $W$ \Term{equivalence}{等价}于给定 $A_U,A_U'$ 之间的这种\Term{duality}{对偶性}。）于是 $W$ 定义\Term{pairing}{配对}
\begin{equation}\tag{14.3.5.4}
T_\ell(A_U)\times T_\ell(A_U')\longrightarrow\mathbb Z_\ell(1),
\end{equation}
并给出\Term{orthogonality}{正交}\Term{theorem}{定理}（7.4.8）：$T_\ell(A_U)^f$ 是 $T_\ell(A_U')^t=T_\ell(T_U')$ 的\Term{orthogonal}{正交补}，其中 $T'$ 为与 $A'$ 关联的\Term{raynaud-extension}{Raynaud 扩张} $G'=A'^\natural$ 的\Term{toric-part}{环面部分}。因此又得到一个\Term{isomorphism}{同构}
\begin{equation}\tag{14.3.5.5}
T_\ell(A_U)/T_\ell(A_U)^f\xrightarrow\sim M_U\otimes_{\mathbb Z}\mathbb Z_\ell,
\end{equation}
其中 $M=D(T')$ 是 $S$ 上与\Term{torus}{环面} $T'$ \Term{dual}{对偶}的\Term{lattice}{格}。故 (14.3.8) 和 (14.3.10) 给出一个\Term{extension}{扩张}结构
\begin{equation}\tag{14.3.5.6}
0\longrightarrow T_\ell(G_U)\longrightarrow T_\ell(A_U)
\longrightarrow M_U\otimes_{\mathbb Z}\mathbb Z_\ell\longrightarrow0.
\end{equation}
现在可以提出如下猜想。

\SourcePage{528}
\begin{SGAresult}{猜想}{14.4}
在 14.3.5 刚阐明的条件下，存在一个\Term{homomorphism}{同态}
\begin{equation}\tag{14.4.1}
i:M_U\longrightarrow G_U,
\end{equation}
使得对每个\Term{prime-number}{素数} $\ell$，\Term{extension}{扩张} (14.3.5.6) 与按照 (14.2.2) 所说明的方式从 $i$ 得到的\Term{extension}{扩张}\Term{isomorphism}{同构}。更确切地说，存在唯一一种方式，能对 14.3.5 所考察类型的每一种情形，赋予一个满足上述条件的\Term{homomorphism}{同态} (14.4.1)，并使 $i$ 的形成与任意\Term{base-change}{基变换} $S'\to S$ 相容；这里的\Term{base-change}{基变换}与\Term{adic-ring}{进}\Term{noetherian}{诺特}环的一个\Term{continuous}{连续}\Term{homomorphism}{同态} $V\to V'$ 关联。
\end{SGAresult}

\textbf{14.4.2。}$i$ 的存在性及其与\Term{base-change}{基变换}的相容性，似乎可由 M. Raynaud 的\Term{rigid-analytic}{刚性解析}理论\Term{pushout}{推出}，因为他正在比 14.1 所陈述的更\Term{generic}{一般}的框架下发展这一理论，也就是像这里一样，考虑\Term{generic}{一般}的\Term{adic-ring}{进}\Term{noetherian}{诺特}环。此外，为证明 14.4 中陈述的存在性和唯一性，立即可归结到只考察 $S_0$ 为 $\mathbb Z$ 上的\Term{finite-type}{有限型}\Term{scheme}{概形}的那些 $S$：事实上，$V$ 是其在 $\mathbb Z$ 上\Term{finite-type}{有限型}的\Term{subring}{子环} $V_i$ 的\Term{filtered-inductive-limit}{滤过归纳极限}；由 Nagata 的结果（EGA IV 7.8），可假设这些环\Term{normal}{正规}。当 $i$ \Term{large-relative-to-cardinal}{足够大}时，情形 $(A,A',W)$ 来自某个 $V_i$ 上的类似情形，或者来自 $V_i$ 对由 $\mathfrak m_i=\mathfrak m\cap V_i$ 所定义的\Term{topology}{拓扑}的\Term{separated-completion}{分离完备化}上的类似情形。这就给出所述的归结。

\SourcePage{529}
\textbf{14.4.3。}现在注意，$S-S_0$ 是一个\Term{jacobson-scheme}{Jacobson 概形}（EGA IV 10.5.8）；所以，对给定情形 $(A,A',W)$，要证明上述问题的两个解 $i,i'$ 相等，只须证明它们在 $U$ 的每个\Term{closed-point}{闭点} $\eta$ 处相等。取 $\eta$ 在 $S$ 中的\Term{closure}{闭包}的\Term{normalization}{正规化} $S'$，则 $S'$ 是\Term{complete}{完备}\Term{trait}{离散赋值环谱}；由于 $S'$ 的\Term{residue-field}{剩余域}为\Term{finite-field}{有限域}，可对 $S$ 上所考察情形在 $S'$ 上的\Term{inverse-image}{逆像}应用 4.3 的唯一性结果。这样就证明了唯一性；在存在性成立的前提下，它蕴涵 $i$ 的形成与\Term{base-change}{基变换}相容。

\textbf{注记 14.5。}尚须以适当的\Term{generic}{一般}性，表述并证明由\Term{rigid-analytic}{刚性解析}方法在 14.4 中引入的\Term{homomorphism}{同态} $i$，与复解析理论中类似的\Term{homomorphism}{同态} (13.5.13) 之间的相容性；更好的做法是找出两者共同的推广！

\SGAsection{}{参考文献}
\begin{description}
\item[{[1]}] M. Artin，《\Term{form}{形式}模\Term{space}{空间}的\Term{algebraization}{代数化} I》（Algebraization of formal moduli I），即将发表。
\item[{[2]}] H. Cartan，《复\Term{space}{空间}族与\Term{analytic-geometry}{解析几何}基础》（Familles d'espaces complexes et fondements de la géométrie analytique），ENS 讨论班，1960/61。
\item[{[3]}] P. Deligne、D. Mumford，《给定\Term{genus}{亏格}的\Term{curve}{曲线}\Term{space}{空间}的\Term{irreducibility}{不可约性}》（The irreducibility of the space of curves of given genus），Pub. Math.，第 36 号（1969）。
\end{description}
\SourcePage{530}
\begin{description}
\item[{[4]}] P. Deligne，关于 \Term{hodge-theory}{Hodge 理论}的筹备中著作（将于 1970 年在 Pub. Math. 发表）。
\item[{[5]}] M. Demazure、J. Giraud、M. Raynaud，《\Term{abelian-scheme}{阿贝尔概形}》（Schémas abéliens），Orsay 讨论班，1967/68，即将发表。
\item[{[6]}] A. Grothendieck，《Brauer 群 III》（Le Groupe de Brauer III），载《\Term{scheme}{概形}\Term{cohomology}{上同调}十讲》（Dix Exposés sur la cohomologie des schémas），North Holland Pub. Cie.
\item[{[7]}] A. Grothendieck，《晶体与\Term{scheme}{概形}的 \Term{de-rham-cohomology}{de Rham 上同调}》（Crystals and the De Rham cohomology of schemes），I. Coates、O. Jussila 整理，载《\Term{scheme}{概形}\Term{cohomology}{上同调}十讲》，North Holland Pub. Cie.
\item[{[8]}] A. Grothendieck，《关于\Term{abelian-scheme}{阿贝尔概形}\Term{homomorphism}{同态}的一个\Term{theorem}{定理}》（Un théorème sur les homomorphismes de schémas abéliens），Inventiones math.，第 2 卷，第 59--78 页（1966）。
\item[{[9]}] P. A. Griffiths，《\Term{hodge-structure}{Hodge 结构}变分报告》（Report on variation of Hodge structures），即将发表。
\item[{[10]}] M. Hakim，《\Term{ringed-topos}{赋环拓扑斯}与\Term{relative-scheme}{相对概形}》（Topos annelés et schémas relatifs），即将刊于 Grundlehren der Math.，Springer（1970）。
\item[{[11]}] J. I. Igusa，《抽象\Term{vanishing-cycles}{消失圈}理论》（Abstract vanishing cycle theory），Proceedings of the Japan Academy，第 34 卷第 9 号（1958），第 589--593 页。
\item[{[12]}] R. Kiehl，《\Term{proper-map}{固有映射}的\Term{finiteness-theorem}{有限性定理}……》（Der Endlichkeitssatz für eigentliche Abbildungen\ldots），Inventiones Math.，第 2 卷，第 191--214 页（1967）。
\item[{[13]}] R. Kiehl。\SGAsourcegap{原 PDF530 文献[13]仅印作者“R. Kiehl”，未印题名及出版信息；保留原有内容，不补建。}
\end{description}
\SourcePage{531}
\begin{description}
\item[{[14]}] S. L. Kleiman，《\Term{algebraic-cycle}{代数圈}与 Weil 猜想》（Algebraic cycles and the Weil conjectures），载《\Term{scheme}{概形}\Term{cohomology}{上同调}十讲》，North Holland Pub. Cie.
\item[{[15]}] S. Lang，《\Term{abelian-variety}{阿贝尔簇}》（Abelian varieties），Interscience Publishers，第 7 号（1959）。
\item[{[16]}] S. Lang、A. Néron，《\Term{function-field}{函数域}上\Term{abelian-variety}{阿贝尔簇}的\Term{rational-point}{有理点}》（Rational points of abelian varieties over function fields），Amer. Jour. Math.，第 81 卷第 1 号，第 95--118 页（1959）。
\item[{[16bis]}] U. I. Manin，《交换\Term{formal-group}{形式群}理论》（Théorie des groupes formels commutatifs），Uspechi Mat. Naouk，第 XVIII 卷，第 4--90 页（1963）。
\item[{[17]}] D. Mumford，《几何\Term{invariants}{不变量}理论》（Geometric invariant theory），Ergebnisse der Mathematik\ldots，第 34 卷，Springer（1965）。
\item[{[18]}] D. Mumford，《抽象 Theta 函数》（Abstract Theta functions），\Term{algebraic-geometry}{代数几何}高级科学讨论班，Bowdoin College，1967（H. Pittie 整理）。
\item[{[19]}] A. Néron，《\Term{local-field}{局部域}与整体域上\Term{abelian-variety}{阿贝尔簇}的\Term{minimal-model}{极小模型}》（Modèles minimaux des variétés abéliennes sur les corps locaux et globaux），Pub. Math.，第 21 号（1964），第 5--128 页。
\item[{[20]}] M. Raynaud，《层的 \Term{euler-poincare-characteristic}{Euler--Poincaré 特征}与\Term{abelian-variety}{阿贝尔簇}的\Term{cohomology}{上同调}》（Caractéristique d'Euler-Poincaré d'un faisceau et cohomologie des variétés abéliennes），载《\Term{scheme}{概形}\Term{cohomology}{上同调}十讲》，North Holland Pub. Cie.
\item[{[21]}] M. Raynaud，《\Term{neron-model}{Néron 模型}》（Modèles de Néron），C. R. Acad. Sci. Paris，第 262 卷，第 413--416 页（1966 年 2 月）。
\item[{[22]}] M. Raynaud，《\Term{picard-functor}{Picard 函子}的\Term{specialization}{特殊化}》（Spécialisation du foncteur de Picard），C. R. Ac. Sc.，第 264 卷，第 941--943 及 1001--1004 页。
\item[{[22bis]}] M. Raynaud，《\Term{picard-functor}{Picard 函子}的\Term{specialization}{特殊化}》（Spécialisation du foncteur de Picard），即将刊于 Pub. Math.
\end{description}
\SourcePage{532}
\begin{description}
\item[{[23]}] M. Raynaud，《\Term{group-scheme}{群概形}与\Term{homogeneous-space}{齐性空间}上的\Term{ample-sheaf}{丰沛层}》（Faisceaux amples sur les schémas en groupes et les espaces homogènes），博士论文，巴黎，1968（将于 1970 年刊于 Springer 的 Lecture Notes）。
\item[{[24]}] M. Raynaud。\SGAsourcegap{原 PDF532 文献[24]仅印作者“M. Raynaud”，未印题名及出版信息；原页\Term{graph}{图像}确认如此，不据同作者其他文献补建。}
\item[{[25]}] M. Rosenlicht，《\Term{algebraic-group}{代数群}的若干基本\Term{theorem}{定理}》（Some basic theorems on algebraic groups），Amer. Jour. Math.，第 78 卷第 2 号，第 401--443 页（1956）。
\item[{[25bis]}] J. P. Serre，《\Term{characteristic}{特征} $p$ 下\Term{abelian-variety}{阿贝尔簇}的若干性质》（Quelques propriétés des variétés abéliennes en car. p），Amer. J. Math.，第 80 卷，1958，第 715--739 页。
\item[{[26]}] J. P. Serre，《\Term{local-field}{局部域}》（Corps locaux），Act. Sci. Ind.，第 1296 号（1962），Herman，巴黎。
\item[{[27]}] J. P. Serre，《\Term{residue-field}{剩余域}\Term{algebraically-closed}{代数闭}的\Term{local-field}{局部域}》（Sur les corps locaux à corps résiduel algébriquement clos），Bull. Soc. Math. France，第 89 卷，1961，第 105--154 页。
\item[{[28]}] J. P. Serre，《阿贝尔 $\ell$-进表示与\Term{elliptic-curve}{椭圆曲线}》（Abelian $\ell$-adic representations and elliptic curves），Benjamin Inc.（1968）。
\item[{[29]}] J. P. Serre、J. Tate，《\Term{abelian-variety}{阿贝尔簇}的\Term{good-reduction}{良约化}》（Good reduction of abelian varieties），Annals of Math.，第 88 卷第 3 号（1968），第 492--517 页。
\item[{[30]}] J. P. Serre、J. Tate，《\Term{elliptic-curve}{椭圆曲线}与\Term{formal-group}{形式群}》（Elliptic curves and formal groups），\Term{algebraic-geometry}{代数几何}暑期研究班笔记，Woods Hole，1964。
\item[{[31]}] J. P. Serre，《泛\Term{morphism}{态射}与 Albanese 簇》（Morphismes universels et variété d'Albanese），载 Chevalley 讨论班，1958/59，第 10 讲。
\item[{[32]}] J. Tate，《$p$-进域上的 WC 群》（WC-groups over p-adic fields），Bourbaki 讨论班，第 156 号，1957。
\item[{[33]}] J. Tate，《$p$-\Term{divisible-group}{可除群}》（$p$-divisible groups），载\Term{local-field}{局部域}会议论文集，Driebergen（1966），Springer。
\item[{[34]}] J. Tate，《\Term{rigid-analytic-space}{刚性解析空间}》（Rigid analytic spaces），IHES 油印本。
\end{description}
